\documentclass[11pt]{article}
\usepackage[a4paper,margin=1in]{geometry}
\usepackage{amsmath,amssymb,amsthm,mathtools}
\usepackage{booktabs}
\usepackage{graphicx}
\usepackage{hyperref}
\usepackage{microtype}

\newtheorem{theorem}{Theorem}
\newcommand{\Cc}{\mathcal C}
\newcommand{\dd}{\,\mathrm d}
\newcommand{\STP}{\mathrm{STP}}
\newcommand{\TP}{\mathrm{TP}}

\title{The gamma-composition kernel is strictly TP$_2$, but not TP$_3$}
\author{Run \texttt{fa\_banach\_001}, agent lane 10}
\date{2026-08-11}

\begin{document}
\maketitle

\section*{Source and open problems}

For \(s,x>0\), write
\[
 Q(s,x)=\frac{\Gamma(s,x)}{\Gamma(s)}
       =\frac1{\Gamma(s)}\int_x^\infty t^{s-1}e^{-t}\dd t
\]
for the regularized upper incomplete gamma function, and let \(Q^{-1}(s,v)\)
denote its inverse in the second variable.  Richards and Uhler
\cite{richards-uhler} introduced, for fixed \(c>0\),
\[
 \Cc_c(u,v)=Q\bigl(u+c,Q^{-1}(u,v)\bigr),
 \qquad u>0,\quad 0\le v<1.
\]
In Remark 4.4 they report substantial numerical evidence that
\((u,v)\mapsto\Cc_c(u,v)\) is both \(\STP_2\) and \(\STP_3\), and pose these as
open problems.  Their decreasing-order convention asks for positive minors
when \(u_1>\cdots>u_r\) and \(v_1>\cdots>v_r\).

\begin{center}
\includegraphics[width=0.96\linewidth]{figures/open_problem_crop.png}
\end{center}

\section*{Answer and proof intuition}

The two orders behave differently.  Differentiate \(\Cc_c\) in \(v\).  The
Jacobian of the inverse gamma CDF cancels the old gamma density, leaving a
positive factor times the \(c\)-th power of the relevant gamma quantile.
Saunders and Moran's strict quantile-ratio theorem \cite{saunders-moran} says
exactly that these powers form an \(\STP_2\) kernel.  Because all rows of
\(\Cc_c\) start at the common value zero when \(v=0\), integration preserves
strict order-two positivity.

Order three fails.  A simple rational parameter choice gives a determinant
near \(-1.21\times10^{-4}\).  We certify its sign by positive-series
enclosures of the incomplete gamma ratios, monotone bracketing of all inverse
quantiles, and outward-rounded interval arithmetic.

\begin{theorem}[Resolution of the two questions]\label{thm:main}
For every \(c>0\), the kernel \(\Cc_c\) is \(\STP_2\) on
\((0,\infty)\times(0,1)\).  It extends to a \(\TP_2\) kernel on
\((0,\infty)\times[0,1)\), but the extension is not strict at \(v=0\).

The proposed general \(\STP_3\) assertion is false.  Indeed, at \(c=1/8\),
the decreasing triples
\[
 (u_1,u_2,u_3)=\left(100,\frac{18}{5},\frac6{25}\right),\qquad
 (v_1,v_2,v_3)=\left(\frac{99}{100},\frac{63}{100},\frac3{20}\right)
\]
give
\[
 \det\bigl(\Cc_{1/8}(u_i,v_j)\bigr)_{i,j=1}^3<0.
\]
Thus the two open problems have answers \emph{yes} at order two and
\emph{no} at order three.
\end{theorem}

\section*{The positive \texorpdfstring{\(\STP_2\)}{STP2} theorem}

Let \(F_u\) be the CDF and \(f_u(x)=x^{u-1}e^{-x}/\Gamma(u)\) the density of a
unit-scale gamma variable of shape \(u\).  Put
\[
 q_u(p)=F_u^{-1}(p),\qquad x_u(v)=q_u(1-v).
\]
Then
\[
 \Cc_c(u,v)=1-F_{u+c}(x_u(v)).
\]
Implicit differentiation of \(F_u(x_u(v))=1-v\) gives
\[
 \frac{\partial x_u}{\partial v}=-\frac1{f_u(x_u(v))}.
\]
Consequently
\begin{equation}\label{eq:derivative}
 h_c(u,v):=\frac{\partial}{\partial v}\Cc_c(u,v)
 =\frac{f_{u+c}(x_u(v))}{f_u(x_u(v))}
 =\frac{\Gamma(u)}{\Gamma(u+c)}x_u(v)^c.
\end{equation}

Saunders and Moran proved that, whenever \(0<\beta<\alpha<1\),
\begin{equation}\label{eq:SM}
 u\longmapsto\frac{q_u(\alpha)}{q_u(\beta)}
 \quad\text{is strictly decreasing on }(0,\infty).
\end{equation}
Fix \(u_1>u_2>0\) and \(v_1>v_2\).  Applying \eqref{eq:SM} with
\(\beta=1-v_1\) and \(\alpha=1-v_2\), and rearranging, yields
\begin{equation}\label{eq:ratio-x}
 \frac{x_{u_1}(v_1)}{x_{u_2}(v_1)}
 >\frac{x_{u_1}(v_2)}{x_{u_2}(v_2)}.
\end{equation}
By \eqref{eq:derivative}, the ratio
\[
 \frac{h_c(u_1,v)}{h_c(u_2,v)}
 =\frac{\Gamma(u_1)\Gamma(u_2+c)}
        {\Gamma(u_1+c)\Gamma(u_2)}
  \left(\frac{x_{u_1}(v)}{x_{u_2}(v)}\right)^c
\]
is therefore strictly increasing in \(v\).  Equivalently, \(h_c\) is
\(\STP_2\).

For completeness, strict positivity survives integration here by an elementary
identity.  Let \(h_1,h_2>0\), suppose \(r=h_1/h_2\) is strictly increasing,
and set \(H_i(v)=\int_0^v h_i(t)\dd t\).  Then
\begin{equation}\label{eq:integral-ratio}
 \left(\frac{H_1}{H_2}\right)'(v)
 =\frac{h_2(v)}{H_2(v)^2}
   \int_0^v\bigl(r(v)-r(t)\bigr)h_2(t)\dd t>0.
\end{equation}
Since \(\Cc_c(u,0)=0\), equation \eqref{eq:derivative} gives
\(\Cc_c(u,v)=\int_0^v h_c(u,t)\dd t\).  Taking
\(h_i(t)=h_c(u_i,t)\) in \eqref{eq:integral-ratio} proves that
\(\Cc_c(u_1,v)/\Cc_c(u_2,v)\) is strictly increasing.  This is precisely the
positivity of every ordered \(2\times2\) minor, proving \(\STP_2\) on the
interior.

If one of the ordered \(v\)'s equals zero, the corresponding column is zero.
The determinant is then nonnegative but not positive.  This proves the exact
boundary statement in Theorem \ref{thm:main}.

\section*{The certified \texorpdfstring{\(3\times3\)}{3x3} counterexample}

At the parameters in Theorem \ref{thm:main}, the kernel matrix is, to 18
decimal places,
{\small
\[
\begin{pmatrix}
0.990341980391842966&0.634744085941809884&0.152888248142902939\\
0.992092691738034760&0.656560587655528763&0.165231530309850930\\
0.999117517495731797&0.785807279598486534&0.233184274842588704
\end{pmatrix}.
\]
}
Its high-precision determinant is approximately
\[
 -0.0001210585212763807470565673499421878.
\]
The following certification removes any reliance on rounded displayed values.
For \(s,x>0\), use the positive series
\begin{equation}\label{eq:series}
 P(s,x):=F_s(x)
 =\frac{x^se^{-x}}{\Gamma(s+1)}
   \sum_{n=0}^{\infty}\frac{x^n}{(s+1)_n}.
\end{equation}
After the ratio of successive terms becomes less than \(1/2\), the remaining
tail is at most twice the next term.  The accompanying script evaluates
\eqref{eq:series} with outward-rounded intervals.  For each pair \((u,v)\), it
first finds numbers \(x_-<x_+\) and certifies
\[
 P(u,x_-)<1-v<P(u,x_+).
\]
Monotonicity then encloses the kernel entry by
\[
 1-P(u+c,x_+)\le \Cc_c(u,v)\le 1-P(u+c,x_-).
\]
Set \(D=\det(\Cc_{1/8}(u_i,v_j))\).  Interval evaluation of its six-term
expansion gives the outward-rounded enclosure
\[
\begin{aligned}
 -1.210585212763807470565673499423\times10^{-4}&\le D,\\
 D&\le-1.210585212763807470565673499421\times10^{-4}.
\end{aligned}
\]
The upper endpoint is strictly negative, proving failure of \(\TP_3\), hence
of \(\STP_3\).

\section*{Verification, novelty, and scope}

The interval certificate is in \texttt{code/interval\_certificate.py} and uses
80 decimal digits of outward-rounded interval arithmetic.  Its approximate
root finder is not trusted: every root is independently bracketed by the
certified series before it enters the determinant calculation.

On 2026-08-11, exact-title, exact-kernel, STP terminology, and citation searches
found the source and its published version but no later paper claiming a
resolution of these two questions.  The run's four lightweight indexes and
local arXiv corpus likewise contained no matching result.  This is a bounded
novelty check, not an exhaustive review of all non-arXiv literature.

The theorem proves \(\STP_2\) for every \(c>0\) and refutes the proposed general
\(\STP_3\) statement with one certified \(c\).  It does not classify those
individual values of \(c\), if any, for which the order-three kernel might
still be strictly totally positive.  Additional searches beyond the certified
small-\(c\) witness produced only cancellation-scale signals for larger \(c\),
so no such classification is asserted.

\section*{Human review recommendation}

Review as a likely-valid full resolution.  The two highest-value independent
checks are: (i) verify the direction change from \eqref{eq:SM} to
\eqref{eq:ratio-x}; and (ii) run the interval certificate and confirm that its
determinant upper endpoint is negative.

\begin{thebibliography}{9}
\bibitem{richards-uhler}
Donald St. P. Richards and Caroline Uhler,
\emph{Loading Monotonicity of Weighted Premiums, and Total Positivity
Properties of Weight Functions}, arXiv:1806.07957 (2018).

\bibitem{saunders-moran}
I. W. Saunders and P. A. P. Moran,
\emph{On the quantiles of the gamma and F distributions},
Journal of Applied Probability \textbf{15} (1978), no.~2, 426--432,
doi:10.2307/3213414.
\end{thebibliography}

\end{document}
